% =============================================================================
% MACROS DE NOTAÇÃO — fonte única para HTML (MathJax) e PDF (LaTeX).
% -----------------------------------------------------------------------------
% • PDF: este arquivo entra no preâmbulo (include-in-header em _quarto.yml).
% • HTML: scripts/gerar_estrutura.py converte este arquivo em
%   estilos/mathjax-macros.html, que configura o MathJax em todas as páginas.
% Regras: use apenas \newcommand (nada de \DeclareMathOperator ou \def) e
% documente cada macro nova no apêndice de notação.
% =============================================================================

% ---- Conjuntos numéricos
\newcommand{\R}{\mathbb{R}}
\newcommand{\N}{\mathbb{N}}
\newcommand{\Z}{\mathbb{Z}}
\newcommand{\Q}{\mathbb{Q}}

% ---- Probabilidade e estatística
\newcommand{\E}{\mathbb{E}}
\newcommand{\Prob}{\mathbb{P}}
\newcommand{\Var}{\operatorname{Var}}
\newcommand{\Cov}{\operatorname{Cov}}
\newcommand{\Corr}{\operatorname{Corr}}
\newcommand{\Avar}{\operatorname{Avar}}
\newcommand{\Normal}{\mathcal{N}}
\newcommand{\iid}{\overset{\text{iid}}{\sim}}
\newcommand{\indep}{\perp\!\!\!\perp}
\newcommand{\pto}{\xrightarrow{\;p\;}}
\newcommand{\dto}{\xrightarrow{\;d\;}}
\newcommand{\qcto}{\xrightarrow{\;\text{q.c.}\;}}
\newcommand{\plim}{\operatorname*{plim}}
\newcommand{\Filt}{\mathcal{F}}

% ---- Álgebra linear
\newcommand{\vb}[1]{\boldsymbol{#1}}
\newcommand{\T}{^{\top}}
\newcommand{\inv}{^{-1}}
\newcommand{\tr}{\operatorname{tr}}
\newcommand{\posto}{\operatorname{posto}}
\newcommand{\diag}{\operatorname{diag}}
\newcommand{\vecop}{\operatorname{vec}}
\newcommand{\Id}{\boldsymbol{I}}

% ---- Otimização e cálculo
\newcommand{\argmax}{\operatorname*{arg\,max}}
\newcommand{\argmin}{\operatorname*{arg\,min}}
\newcommand{\dd}{\mathrm{d}}
\newcommand{\pd}[2]{\frac{\partial #1}{\partial #2}}
\newcommand{\defeq}{\mathrel{:=}}

% ---- Séries temporais
\newcommand{\Lag}{\mathrm{L}}

% ---- Diversos
\newcommand{\abs}[1]{\left\lvert #1 \right\rvert}
\newcommand{\norm}[1]{\left\lVert #1 \right\rVert}
\newcommand{\inner}[2]{\left\langle #1, #2 \right\rangle}
\newcommand{\diam}{\operatorname{diam}}
\newcommand{\interior}{\operatorname{int}}
\newcommand{\ind}[1]{\mathbf{1}\{#1\}}

% ---- Econometria (siglas em português)
\newcommand{\SQT}{\mathrm{SQT}}
\newcommand{\SQE}{\mathrm{SQE}}
\newcommand{\SQR}{\mathrm{SQR}}
\newcommand{\ep}{\operatorname{ep}}
